\part{Annexes}

\chapter{Refaire le projet depuis zéro}

Voici l'ordre dans lequel aborder un projet de ce type. Chaque étape produit
quelque chose qui marche : c'est ce qui permet de ne jamais se retrouver avec
mille lignes dont rien ne fonctionne.

\begin{enumerate}
\item \textbf{Un conteneur SSH à la main.} Écrivez un Dockerfile qui installe
  \cmd{openssh-server}, lancez-le avec \cmd{-p 2222:22}, connectez-vous.
  Objectif : comprendre la publication de port.
\item \textbf{Plusieurs services dedans.} Ajoutez \cmd{supervisord} et un
  terminal web. Objectif : comprendre PID 1.
\item \textbf{Une page Flask qui lance ce conteneur.} Un bouton, un
  \cmd{subprocess}, un message. Pas de base, pas de comptes.
\item \textbf{Une base.} Enregistrez ce que vous créez. Séparez dès maintenant
  la machine de la location : ce choix est difficile à rattraper après coup.
\item \textbf{Des comptes.} Inscription, connexion, et surtout le cloisonnement
  --- un utilisateur ne voit que ses machines.
\item \textbf{L'expiration.} Un script séparé qui boucle et détruit. Testez-le
  avec des locations d'une minute.
\item \textbf{Les tests.} Maintenant, pas plus tard : vous avez assez de code
  pour que les régressions commencent.
\item \textbf{Un second nœud.} C'est là que tout change : SSH, ordonnanceur,
  agent. Faites-le avant d'avoir trop de code supposant un hôte unique.
\item \textbf{Ansible.} Automatisez ce que vous faites à la main depuis le
  début. Visez \cmd{changed=0} au second passage dès le premier jour.
\item \textbf{La chaîne d'intégration.} Commencez par lint et tests, ajoutez
  le reste progressivement.
\item \textbf{La haute disponibilité.} Éteignez un nœud et regardez ce qui se
  passe. Corrigez. Recommencez.
\end{enumerate}

\begin{retenir}
L'ordre compte plus qu'on ne croit. Ajouter le multi-nœuds à la fin oblige à
réécrire tout le code qui appelait \cmd{docker} sans se demander « sur quelle
machine ? ». Ajouter les tests à la fin revient à ne pas les ajouter.
\end{retenir}

\chapter{Commandes utiles}

\section{Le projet}

\begin{lstlisting}[language=bash]
# --- Pile complete en local --------------------------------------------------
cp .env.example .env                     # puis renseigner les secrets
docker compose up -d
docker compose --profile build build     # construire les six images
docker compose logs -f web
docker compose down -v                   # tout effacer, volumes compris

# --- Developpement -----------------------------------------------------------
cd 2_webapp
python3 -m venv .venv && .venv/bin/pip install -r requirements-dev.txt
.venv/bin/python app.py                  # serveur de developpement

# --- Controles, les memes que la chaine --------------------------------------
pytest                                   # 169 tests
ruff check .
bandit -c pyproject.toml -r .
pip-audit -r requirements.txt
alembic check
./tests/dast.sh http://127.0.0.1:5000
./tests/integration.sh http://127.0.0.1:8088

# --- Cluster -----------------------------------------------------------------
vagrant up
cd 3_ansible && ansible-playbook site.yml --ask-vault-pass
ansible-playbook site.yml --ask-vault-pass    # doit donner changed=0
vagrant halt worker1                     # le test de panne
\end{lstlisting}

\section{Diagnostiquer}

\begin{lstlisting}[language=bash]
# Que fait reellement un conteneur ?
docker inspect <conteneur> | jq '.[0].HostConfig | {Memory, NanoCpus, CapDrop, PidsLimit}'

# Pourquoi ne demarre-t-il pas ?
docker logs <conteneur>

# Quelles machines la plateforme reconnait-elle comme siennes ?
docker ps -a --filter "label=insacloud.role=rented-machine"

# Les services sur le controleur
systemctl status insacloud insacloud-faucheur nginx
journalctl -u insacloud-faucheur -f

# L'etat du parc, vu de l'exterieur
curl -s http://127.0.0.1:8088/health | jq
curl -s -H 'Accept: application/json' http://127.0.0.1:8088/workers | jq
\end{lstlisting}

\chapter{Erreurs fréquentes}

\begin{center}
\small
\begin{tabularx}{\textwidth}{@{}XX@{}}
\toprule
\textbf{Symptôme} & \textbf{Cause probable} \\
\midrule
\cmd{exec format error} ou \cmd{no such file or directory} sur un script qui
existe pourtant &
Le script a des fins de ligne Windows (CRLF). Le \cmd{\textbackslash r} fait
partie du nom de l'interpréteur dans le shebang. Corriger avec un
\cmd{.gitattributes} : \cmd{* text=auto eol=lf} \\
\addlinespace
\cmd{port is already allocated} &
Un conteneur précédent n'a pas été nettoyé, ou deux requêtes ont choisi le même
port. L'index unique partiel règle le second cas \\
\addlinespace
\cmd{table users already exists} au démarrage &
Plusieurs processus Gunicorn créent le schéma en même temps. Utiliser
\cmd{--preload} et tolérer l'erreur \\
\addlinespace
Le conteneur démarre puis s'arrête aussitôt &
Son processus principal s'est terminé. Un conteneur vit tant que son PID 1 vit \\
\addlinespace
\cmd{permission denied} sur \cmd{/var/run/docker.sock} &
Le compte n'est pas dans le groupe \cmd{docker}. En conteneur, utiliser
\cmd{group\_add} avec le GID de l'hôte \\
\addlinespace
Ansible affiche \cmd{changed} à chaque passage &
Un module \cmd{command} ou \cmd{shell} sans \cmd{creates:} ni
\cmd{changed\_when:} \\
\addlinespace
Le site répond en HTTP mais pas en HTTPS &
Certificat absent ou illisible par nginx, ou port 443 fermé par UFW \\
\addlinespace
Tout le monde est bloqué par l'anti-force-brute &
\cmd{ProxyFix} n'est pas actif : toutes les requêtes semblent venir du proxy \\
\addlinespace
Les tests passent ensemble mais échouent seuls (ou l'inverse) &
Un test laisse des données derrière lui. Remettre la base à zéro avec une
fixture \cmd{autouse} \\
\addlinespace
La connexion SSH échoue juste après la location &
Docker publie le port avant que \cmd{sshd} écoute. Réessayer avec une limite,
plutôt qu'attendre une durée fixe \\
\bottomrule
\end{tabularx}
\end{center}

\chapter{Conformité au cahier des charges}

\begin{center}
\small
\begin{tabularx}{\textwidth}{@{}llX@{}}
\toprule
\textbf{\#} & \textbf{Fonctionnalité} & \textbf{Réalisation} \\
\midrule
1 & Préparation et environnement & dépôt Git, environnement virtuel, dépendances épinglées \\
2 & API Web Flask & modules séparés, comptes, sessions, tableau de bord \\
3 & Base de données & 7 tables en SQLAlchemy, migrations Alembic, index métier \\
4 & Conteneurisation Docker & 4 Dockerfiles, 6 images, durcissement complet \\
5 & Infrastructure Vagrant & 4 VM, réseau privé, 4 providers \\
6 & Automatisation Ansible & 9 rôles, 4 playbooks, Vault, idempotence mesurée \\
7 & Worker Agent & inscription spontanée, battement de cœur, 3 états \\
8 & Ordonnanceur & nœud joignable le moins chargé, départage déterministe \\
9 & Service de location & durée 1--120 min, prolongation, arrêt \\
10 & Déploiement des instances & la commande imposée, par \cmd{subprocess} \\
11 & Accès SSH & par mot de passe ou par clé, plus terminal web et bureau \\
12 & Expiration & le Faucheur, toutes les 10 s, avec réconciliation \\
13 & Haute disponibilité & 3 niveaux, scénario de l'énoncé vérifié \\
14 & Sécurité DevSecOps & SAST, SCA, secrets, DAST, scan d'image, SBOM \\
15 & CI/CD & 11 travaux, publication conditionnelle, déploiement vérifié \\
16 & Tests d'intégration & 169 tests unitaires + 6 scénarios de bout en bout \\
17 & Documentation & README, guide d'API, schéma de base, ce guide \\
\bottomrule
\end{tabularx}
\end{center}

\section*{Au-delà de l'énoncé}

\begin{itemize}
\item Architecture distribuée à trois workers, alors qu'un seul hôte suffisait.
\item Terminal web et bureau graphique complet dans le navigateur, en plus du
  SSH demandé.
\item Coffre à secrets : mot de passe root chiffré, réaffichable après
  re-authentification.
\item Watchdog systemd d'auto-réparation du service.
\item Inventaire dynamique interrogeant Vagrant.
\item Sauvegarde du travail avant destruction, téléchargeable sept jours.
\item Schémas engendrés depuis leur description, et contrôlés par la chaîne.
\end{itemize}

\chapter{Glossaire}

\begin{description}
\item[Agent] petit programme tournant sur chaque worker, qui l'inscrit auprès
  du contrôleur et signale régulièrement sa présence.
\item[Alembic] outil de migration de schéma pour SQLAlchemy.
\item[Ansible] outil de configuration qui applique un état décrit en YAML, par
  SSH, de façon idempotente.
\item[Capacité (\emph{capability})] fraction des pouvoirs de root, attribuable
  indépendamment des autres.
\item[cgroup] mécanisme du noyau Linux limitant les ressources qu'un groupe de
  processus peut consommer.
\item[CSP] \emph{Content Security Policy} : en-tête qui restreint les sources
  de scripts, styles et images qu'une page peut charger.
\item[CSRF] attaque où un site tiers fait exécuter une action au nom d'un
  utilisateur connecté ailleurs.
\item[CVE] identifiant mondial d'une faille de sécurité connue.
\item[DAST] analyse de sécurité de l'application \emph{en fonctionnement}.
\item[Faucheur] le démon du projet qui détruit les machines expirées,
  réconcilie Docker et la base, et surveille le parc.
\item[Fernet] recette de chiffrement symétrique authentifié de la bibliothèque
  \cmd{cryptography}.
\item[Gunicorn] serveur WSGI multi-processus qui exécute l'application Flask en
  production.
\item[Hachage] transformation à sens unique d'un mot de passe, non réversible.
\item[Idempotence] propriété d'une opération dont la deuxième application ne
  change rien de plus que la première.
\item[Image] modèle figé en lecture seule, à partir duquel on crée des
  conteneurs.
\item[Namespace] mécanisme du noyau Linux qui cloisonne ce qu'un processus peut
  voir.
\item[ORM] bibliothèque faisant correspondre classes Python et tables SQL.
\item[Proxy inverse] serveur placé devant l'application, qui reçoit les
  requêtes à sa place.
\item[SAST] analyse de sécurité du code source, sans l'exécuter.
\item[SBOM] inventaire exhaustif des composants d'un logiciel livré.
\item[SCA] confrontation des dépendances aux bases de vulnérabilités connues.
\item[supervisord] gestionnaire de processus utilisé comme PID 1 dans les
  machines louées.
\item[Vagrant] outil de création de machines virtuelles décrites dans un
  fichier.
\item[Vault] fichier de secrets chiffré par Ansible, versionnable sans danger.
\item[WSGI] contrat d'interface entre un serveur Python et une application
  Python.
\end{description}

\chapter{Pour aller plus loin}

\section{Ce que ce projet ne fait pas}

\begin{description}
\item[Pas de stockage partagé] la reprise après panne ne rend pas les fichiers.
  Un volume répliqué (NFS, Ceph, Longhorn) le permettrait, au prix d'une
  complexité bien supérieure.
\item[Le contrôleur reste un point unique] s'il tombe, le tableau de bord
  disparaît et les locations cessent d'expirer. Les machines louées, elles,
  continuent de tourner. Une vraie haute disponibilité demanderait plusieurs
  contrôleurs et une base répliquée.
\item[Certificat auto-signé] remplaçable par \cmd{mkcert} ou une autorité
  interne ; Let's Encrypt exigerait un nom de domaine public.
\item[Pas d'authentification à deux facteurs], ni de quotas cumulés par
  utilisateur, ni de récupération de mot de passe.
\item[Docker publie ses ports directement dans iptables] : un filtrage strict
  des conteneurs passerait par la chaîne \cmd{DOCKER-USER}.
\end{description}

\section{Les suites naturelles}

\begin{itemize}
\item Remplacer l'ordonnanceur maison par Kubernetes, et comprendre au passage
  ce que Kubernetes résout vraiment.
\item Ajouter de la métrologie : Prometheus pour les mesures, Grafana pour les
  courbes, et des alertes sur la saturation du parc.
\item Centraliser les journaux des quatre nœuds, pour pouvoir enquêter sans se
  connecter machine par machine.
\item Facturer à la minute, ce qui suppose de fiabiliser la mesure du temps
  réellement consommé.
\end{itemize}

\section{Le dépôt}

\begin{center}
\url{https://github.com/aziz3r/insacloud}
\end{center}

\vspace{1cm}
\begin{center}
\color{gristexte}\small
Ce guide décrit l'état du projet au \today.\\
Les extraits de code sont tirés du dépôt sans retouche.
\end{center}
